\begin{statement}{E.47}{Proved}{The complete three-adic class-three correction-image fibre and its exact exponent-three image}
Let
\[
 \begin{gathered}
 C_3=H^{(3)}/\overline{\gamma_4(H^{(3)})},\qquad
 K_3=\ker(\widehat H\longrightarrow C_3),\\
 \mathcal B_3=\widehat B/K_3,\qquad
 \mathscr J_3^{\infty}
 =\operatorname{im}(\widehat B\longrightarrow
                  \mathscr D\times\mathcal B_3).
 \end{gathered}
 \tag{E.47.1}
\]
Here $H^{(3)}$ is the maximal pro-three quotient,
$\mathscr D$ is the actual E.24 common source, and
$H\triangleleft F\triangleleft B$ and all based maps
are those of E.26.  Fix any one of E.24's $432$
multiplier-one source values, retaining its literal
commutator word.

The complete set of boundary correction images in
$\mathscr J_3^{\infty}$ satisfying the two order
relations has an explicit continuous parametrization
\[
                         \mathbb Z_3^{70}.
 \tag{E.47.2}
\]
Its projection to the full E.37 class-two fibre is
given by the four exact equations
\textup{(E.47.10)} and the two residue equations
\textup{(E.47.15)} below.  Every point of this
projection has a complete central fibre affine
under a free $\mathbb Z_3$ module of rank $64$.

The image in E.44's actual exponent-three target
has exactly $3^{25}$ elements per source value.
Its complete criterion consists of the E.37
reduction conditions, E.44's equations, and the
two additional equations \textup{(E.47.22)}.
All $24$ retained CM rows, with all $18$ source
points in each row, survive.
There is a specified E.44 shadow, compatible
with the entire class-two tower, which has no
lift to this full class-three tower:
\[
 t=0,\qquad Y_{41}=1,\qquad
 Y_j=0\quad(37\leq j\leq44,\ j\ne41).
 \tag{E.47.3}
\]
Its exact finite central value is printed in
\textup{(E.47.24)}.  The failure persists after
changing every higher class-two digit.

All the stated correction images are generating
in the actual common finite sources.  They are
realized by single profinite commutator words,
simultaneously with E.41's full two-adic class-two
family, E.42's full pro-$p$ factors for $p\geq5$,
the E.43 refinement, and the five-point Hurwitz
factor through the exact maps of E.46.26.
\end{statement}
\begin{proof}
We first give the actual class-three group and
the exact source comparison.  Let $V=\mathbb Z_3^{10}$
with ordered basis $e_1,\ldots,e_{10}$ and let
$L_3(V)$ be the ordinary degree-three free Lie
lattice.  Its $330$ ordered Lyndon brackets are
expanded in the $1000$ tensor words of length
three; each has leading coefficient one.
Use scaled logarithmic coordinates
\[
 \log g=(v,Z/2,T/12),\qquad
 v\in V,\quad Z\in\Lambda^2V,\quad T\in L_3(V).
 \tag{E.47.4}
\]
Multiplication is the following integral polynomial:
\[
 \begin{split}
 (v,Z,T)(w,W,U)=\bigl(&v+w,\ Z+W+v\wedge w,\\
 &T+U+3([v,W]-[w,Z])+[v-w,v\wedge w]\bigr).
 \end{split}
 \tag{E.47.5}
\]
The inverse negates all three coordinates.
This formula follows directly by collecting the
degree-one, degree-two and degree-three terms of
the Baker--Campbell--Hausdorff identity; the
displayed factors retain its denominator three.

Here is the precise group lattice in these
coordinates.  Write the product of the ordered
generator powers $\prod_i\eta_i^{v_i}$ as
$(v,Z_0(v),T_b(v))$, so $(Z_0)_{ij}=v_iv_j$.
Define the integral linear map
\[
 \begin{gathered}
 \mathcal C(e_i\wedge e_j)
 =[e_i,[e_i,e_j]]+[e_j,[e_i,e_j]],
 \\
 T_0(v,Z)=T_b(v)+3\mathcal C(Z-Z_0)+3[v,Z-Z_0].
 \end{gathered}
 \tag{E.47.6}
\]
The actual Hall normal form gives exactly
\[
 C_3=\{(v,Z,T):T-T_0(v,Z)\in3L_3(V)\},
 \qquad
 \xi=\frac{T-T_0(v,Z)}{12}\in L_3(V).
 \tag{E.47.7}
\]
Indeed the ordered commutator exponents are
$(Z-Z_0)/2$; their logarithms and their products
with $\prod_i\eta_i^{v_i}$ give
\textup{(E.47.6)}.  The remaining ordinary Hall
commutators of weight three have actual logarithms
equal to the Lyndon basis, and hence scaled $T$
coordinates equal to $12$ times that basis.
Since $2$ and $4$
are units in $\mathbb Z_3$, these expressions
give inverse continuous coordinate maps between
the Hall normal form and \textup{(E.47.7)}.
They prove the comparison with the actual pro-three
quotient, including the asserted congruence.
Reduction of the integral coordinates at increasing
three-power precision gives a cofinal system of
finite continuous images of this group.

The based actions $S,R$ are obtained by substituting
the literal words
$x\mapsto y,y\mapsto x$ and
$x\mapsto y,y\mapsto(xy)^{-1}$ in all ten
Schreier words.  Their full $385$ by $385$ matrices
in \textup{(E.47.4)} are integral and satisfy
$S^2=R^3=1$ over $\mathbb Z$.
Their first $55$ rows are exactly the E.37
doubled-log action.  The last diagonal blocks
$S_3,R_3$ are the ordinary free Lie actions on
$L_3(V)$; all off-diagonal blocks from the
literal substitutions are retained.

The source image is also exact.  The $45$ adjusted
relators $u_j\in H'$ of E.33 have central images
forming a three-adic basis.  The $450$ literal
words
\[
                  [\eta_i,u_j^4]
       \in\ker(B\longrightarrow\mathscr D)
                       \cap\gamma_3(H)
 \tag{E.47.8}
\]
map to a spanning set of $L_3(V)$ modulo three.
The retained bracket matrix has rank $330$ and
specifies a unit $330$-column minor.  This rank
also follows from the free Lie brackets
$[V,\Lambda^2V]=L_3(V)$ and the unit degree-two
matrix of the $u_j$.  The fourth powers clear
the binary source image and preserve its ternary
kernel.  Their three-supported continuous powers
therefore supply every element of $L_3(V)$ while
retaining all class-two factors and all pro-$p$
factors for $p\ne3$.
Together with the exact E.37 source image, this
proves that the new source image is its full
inverse image under $C_3\to C_3/\gamma_3(C_3)$.
Thus the only affine source restriction is the
one already proved in E.37 on $v$; there is no
uncomputed central source condition.

We solve the two order equations with this
entire source image.  Over an E.37 point $(v,Z)$
put $g_0=(v,Z,T_0(v,Z))$, and let $E_s,E_r$ be
the scaled third coordinates of
$(Sg_0)g_0$ and
$g_0^{-1}R(g_0^{-1})R^2(g_0^{-1})$.
Their lower coordinates vanish.  The equations
for $g=(v,Z,T_0+12\xi)$ are
\[
 M\xi=b,\qquad
 M=\begin{pmatrix}1+S_3\\-1-R_3-R_3^2\end{pmatrix},
 \qquad b=-\frac1{12}\binom{E_s}{E_r}.
 \tag{E.47.9}
\]
The exact $660$ by $330$ matrix has Smith factors
$1^{264},3,3,0^{64}$.  In particular its rational
rank is $266$, its integral kernel has rank $64$,
and its image has index $9$ in its three-saturated
rational image.
The errors lie in
$L_3(V)^{S_3=1}\oplus L_3(V)^{R_3=1}$:
cyclic permutation of each central norm word
gives the same word.  The two invariant dimensions
are $164$ and $106$; the traces are $-2$ and $-6$.
The common dual invariant space has rank four.
For its four saved primitive rows $\Lambda$, the
identities
$(3\Lambda,2\Lambda)M=0$ give four independent
conditions on this $270$-dimensional error space.
They therefore characterize rational membership
in the image of $M$.

We give those four conditions as small explicit
equations on the full E.37 chart.  Let $c\in\{0,1,2\}$
be the actual affine source residue computed in
E.37 and put
\[
 U=t_1+t_2=c+3w,\quad u=t_2,\quad t_3=u+3e,
 \qquad \Phi=e+U^2+3Ue+9e^2=0.
 \tag{E.47.10a}
\]
For every $U\in\mathbb Z_3$ there is a unique
$e\in\mathbb Z_3$ satisfying this equation:
$e\equiv-U^2\pmod3$ and
$\partial\Phi/\partial e\equiv1\pmod3$.
The substitution sends the E.37 equation
$t_3-t_2+3Q(t)=0$ exactly to $3\Phi=0$.
Set $a=(Y_{37},\ldots,Y_{44})^t$.
The complete rational criterion is
\[
                         A(U,e)a+h(U,u,e)=0,
       \qquad A=A_0+U A_U+e A_e,
 \tag{E.47.10}
\]
where
\[
 A_0=\begin{pmatrix}
 -1&-1&-2&-1&-1&0&0&1\\
 3&3&6&3&5&0&0&-1\\
 -7&-7&-10&-7&-13&2&0&-1\\
 -4&-6&-8&-4&-6&-2&2&0
 \end{pmatrix},\qquad
 A_U=\begin{pmatrix}
 0&0&0&0&0&0&0&0\\0&0&0&0&0&0&0&0\\
 16&24&8&12&0&-4&-10&-8\\
 -30&-36&-12&-24&0&12&12&24
 \end{pmatrix},
 \tag{E.47.11}
\]
\[
 A_e=\begin{pmatrix}
 0&0&0&0&0&0&0&0\\0&0&0&0&0&0&0&0\\
 6&36&12&0&0&12&-24&24\\
 -54&-108&-36&-36&0&0&54&0
 \end{pmatrix}.
 \tag{E.47.12}
\]
The four constant polynomials, reduced by the
displayed monic relation $\Phi$ in the variable
$U$, are
\[
\begin{aligned}
 h_1={}&-U+6e^2+6eu+2u^2+u,\\
 h_2={}&U(12e+8u+1)-6e^2-30eu-14u^2-u,\\
 h_3={}&U(12e^2-24eu-150e-14u^2-76u-1)\\
       &-306e^3-360e^2u-372e^2-102eu^2
                       +12eu-28e+38u^2+u,\\
 h_4={}&2\bigl[U(-63e^2-36eu+39e-3u^2+31u-1)\\
       &\hspace{8mm}+135e^3+216e^2u+333e^2+72eu^2
                       +120eu+30e+12u^2+3u\bigr].
\end{aligned}
 \tag{E.47.13}
\]
These equations are obtained by applying the
four common dual rows to the exact integral
norms, substituting the E.37 Smith origin, and
performing the saved invertible rational row
change.  All row divisions and polynomial
multiples of $\Phi$ are retained.  This gives
an equality of zero loci over $\mathbb Q_3$,
not a discarded integral condition.

Number the columns of $A$ from zero.  For
$P=(0,1,2,4)$ and $F=(3,5,6,7)$, its pivot
minor satisfies the exact identity
\[
 \det A_P=-16\bigl(30\Phi+1-3U-3e\bigr).
 \tag{E.47.14}
\]
It is a three-adic unit at every point of the
chart.  Thus $a_F=b=(b_0,b_1,b_2,b_3)^t$ is
free and $a_P=-A_P^{-1}(h+A_Fb)$ gives every
rationally soluble integral point.

The two remaining integral Smith tests, after
this substitution, are exactly
\[
 \begin{pmatrix}2&0&1&1\\c+2&1&-c&-c\end{pmatrix}
                 \bar b=\rho(c,\bar u)
                       \quad\hbox{in }\mathbb F_3^2,
 \qquad
 \begin{array}{c|ccc}
 &\bar u=0&\bar u=1&\bar u=2\\\hline
 c=0&(0,0)&(1,1)&(0,2)\\
 c=1&(1,0)&(0,0)&(0,2)\\
 c=2&(2,0)&(2,0)&(0,1)
 \end{array}.
 \tag{E.47.15}
\]
For completeness, the two functionals are
constructed on the three-saturated image of
$M$, rather than extended from a chosen rational
complement.  Starting with $266$ independent
columns of $M$, two recorded divisions by three
give a matrix $N$ whose reduction has rank $266$.
Hence its columns are a basis of
$(\operatorname{im}M_{\mathbb Q_3})\cap\mathbb Z_3^{660}$.
Choose the recorded $266$ rows $I$ for which $N_I$
is invertible over $\mathbb Z_3$.  The annihilator of the image
of $N_I^{-1}M_I$ modulo three has dimension two.
Pulling its rows back by $N_I^{-1}$ gives the
two saved tests.  The Smith factors show that
their vanishing is both necessary and sufficient
for integral membership.
In their polynomial evaluation the substitution
$t_3=u+3e$ is made before division by three.
The resulting two weighted error polynomials
have integral coefficients after that division.
Their residues depend only on $(c,\bar u,\bar b)$;
the exact coefficient matrices give
\textup{(E.47.15)} in all nine charts.

Every chart is soluble, since the first two
columns of the matrix in \textup{(E.47.15)}
have determinant two.  More explicitly choose
$u,w,b_2,b_3,\ell_0,\ell_1\in\mathbb Z_3$.
Use the digits $d_0,d_1\in\{0,1,2\}$ defined by
\[
 \begin{gathered}
 d_0=2(\rho_1-\bar b_2-\bar b_3),\qquad
 d_1=\rho_2-(c+2)d_0+c(\bar b_2+\bar b_3)
                                  \quad\hbox{in }\mathbb F_3,\\
 b_0=d_0+3\ell_0,\qquad b_1=d_1+3\ell_1.
 \end{gathered}
 \tag{E.47.16}
\]
Equations \textup{(E.47.10a), (E.47.14), (E.47.16)}
give a continuous bijection from $\mathbb Z_3^6$
onto the entire projected E.37 locus.  Its inverse
recovers $u,w$, the four entries of $a_F$, and
divides the last two digit discrepancies by three.

We next supply the central solution, including
its integral right inverse.  The same two
saturation operations give an integral matrix
$P_0$ with
\[
                        M P_0=3N.
 \tag{E.47.17}
\]
The saved numerator initially has denominator
nine; every numerator entry is divisible by
three, giving exactly \textup{(E.47.17)}.
Each saturation step enlarges the selected
column lattice by index three: the coefficient
of the replaced column in the saved relation
is a three-adic unit.  Thus the original
selected-column lattice has index nine in
$N\mathbb Z_3^{266}$.  The full image of $M$
has that same index by its Smith factors and
contains the selected columns.  The two
lattices are therefore equal.  This proves
that $P_0/3$, on the specified $N$ coordinates,
is an integral right inverse exactly on the
full image of $M$.
For the right side $b$ in \textup{(E.47.9)} set
\[
 \xi_0=\frac13P_0N_I^{-1}b_I,\qquad
 \xi=\xi_0+Kz,\quad z\in\mathbb Z_3^{64},
 \qquad T=T_0+12\xi.
 \tag{E.47.18}
\]
Here $K$ is the saved integral kernel matrix.
It satisfies $MK=0$ and has a unit $64$-row
minor.  Its rational columns span $\ker M$, so
these properties prove that its columns form
the complete three-adic kernel lattice.
On the rational image, $NN_I^{-1}b_I=b$.
The tests \textup{(E.47.15)} state exactly that
the numerator of $\xi_0$ is divisible by three.
Consequently \textup{(E.47.18)} is an actual
integral solution of \textup{(E.47.9)}, and
subtracting it from any solution gives a unique
$Kz$.  This proves the $64$-dimensional central
fibre and the complete $70$-parameter
parametrization \textup{(E.47.2)}.

We now determine the precise finite projection.
The map to E.44 is the actual Hall/Magnus map
\[
 \begin{gathered}
 (v,Z,T)\longmapsto
 \left(\bar v,\,2\bar Z,
       t_{\rm can}(\bar v,2\bar Z)+\Pi\bar\xi\right),
 \\
 \Pi(L)_{ijk}=[X_iX_jX_k]L-[X_jX_iX_k]L
                  \quad(i<j<k).
 \end{gathered}
 \tag{E.47.19}
\]
It is obtained by expanding the literal Hall
normal form.  The tensor-coordinate definition
is the cube-span quotient proved in E.44, so
it respects multiplication and both based
actions.  The matrix $\Pi K$ has rank $21$
over $\mathbb F_3$; E.44's central norm kernel
has rank $22$.  Thus over each fixed exact
E.37 point in the locus just proved, the
finite central image is the explicit affine
space
\[
      t_{\rm can}+\Pi\bar\xi_0+\operatorname{im}(\Pi K),
 \tag{E.47.20}
\]
of size $3^{21}$.  This is the complete
fixed-lower-point comparison.

Changing a higher class-two digit supplies the
remaining direction when only the finite E.44
point is prescribed.  The central error is
affine in $Z$; write its coefficient after
division by three as $C(t)$.  The exact rational
graph has $Z$-derivative
$D=V_M{}_{[:,37:44]}\,\partial a/\partial b$.
Changing $b$ to $b+3d$ keeps all integral tests
and all class-two residues.  Formula
\textup{(E.47.18)} gives its central derivative
modulo three as
\[
                   -P_0N_I^{-1}(C(t)D)_I.
 \tag{E.47.21}
\]
The factor $1/4$ from \textup{(E.47.9)} has
residue one.  A saved linear functional that
annihilates $\Pi K$ evaluates the four projected
columns of \textup{(E.47.21)} as $(2,0,1,1)$
in each of the nine charts.  In particular
changing $\ell_0$ by one supplies a complement
to $\operatorname{im}(\Pi K)$.  The full finite
central kernel is therefore obtained, using
this actual change of lower parameters and
then the actual central kernel corrections.

It remains to identify which finite lower
points occur.  On E.44's locus compatible
with E.37, the four rational residue equations
and the two integral tests add exactly two
independent equations.  With
$a=(\bar Y_{37},\ldots,\bar Y_{44})^t$, they are
\[
 \begin{pmatrix}
 2&2&1&2&2&0&0&1\\
 c+2&2&2c+2&2&2&2c+2&2c&c+2
 \end{pmatrix}a
 =-\binom{2c+2\bar u^2+\bar u}
 {(c+2)\bar u^2+(2c+1)\bar u+c^2+2c}
                   \quad\hbox{in }\mathbb F_3^2.
 \tag{E.47.22}
\]
These are rows one and three of
\textup{(E.47.10)} reduced modulo three.
For each of the three values of $\bar u$,
their joint rank with E.44.14 is six, and the
joint affine system has nine solutions for
$a$.  Formula \textup{(E.47.16)} realizes all
of them.  There are consequently $3^3$ lower
images per source, each with the entire
$3^{22}$ finite central fibre by
\textup{(E.47.21)}.  This proves the exact
count $3^{25}$ and the complete image criterion.
For comparison, E.44 has $3^{28}$ points per
source, of which $3^{27}$ are compatible with
the full E.37 tower.

The new nonlifting example has the integral
E.37 central coordinate
\[
 Z_*=-e_1\wedge e_7-e_2\wedge e_7-e_3\wedge e_6
       -e_3\wedge e_{10}-e_4\wedge e_7-e_5\wedge e_7
       +e_7\wedge e_8+e_7\wedge e_9.
 \tag{E.47.23}
\]
It is the column $V_Me_{41}$, so it has exactly
the free Smith coordinates in
\textup{(E.47.3)}.  Use the actual source
$(\text{row }0;\,\text{ternary }(0,0);\,\text{binary }0)$,
whose affine residue is zero.
In the E.44 coordinates the finite point is
$(0,2\bar Z_*,t_*)$, where the following array
lists $t_*$ consecutively in the lexicographic
order of the triples $1\leq i<j<k\leq10$:
\[
\begin{pmatrix}
0&2&2&0&2&2&0&2&2&0\\
1&2&2&0&1&1&0&1&2&2\\
1&1&1&1&0&2&1&0&2&0\\
1&1&0&0&2&0&1&0&0&1\\
0&0&1&1&2&2&1&0&0&2\\
2&0&2&2&0&0&0&0&2&0\\
0&0&0&0&2&1&1&0&0&0\\
0&2&0&0&2&1&0&0&0&2\\
1&1&0&0&0&0&0&1&1&1\\
0&0&0&0&2&0&0&0&0&0\\
0&0&0&0&0&0&0&0&0&0\\
0&0&0&0&0&0&0&0&0&0
\end{pmatrix}.
 \tag{E.47.24}
\]
Both E.44 norms vanish exactly.  Its full
class-two tower point satisfies E.44.14,
so E.44's literal central corrections realize
it in the actual common source, retaining
that whole tower.  In \textup{(E.47.22)} its
left side is $(2,2)^t$ and its right side is
zero.  Since these are conditions on the
finite lower coordinates, no change of any
higher digit can give a full class-three lift.
This proves the stated failure for an actual
peripheral generating shadow.

We finish with the simultaneous source and
generation assertions.  The new source-kernel
equality proved using \textup{(E.47.8)} lifts
every prescribed E.37 point to every permitted
class-three central value.  The corresponding
three-supported powers preserve the whole
two-adic class-two family and all full pro-$p$
factors with $p\ne3$.  Compactness of the closed
source-kernel image, or the explicit products
of the $330$ basis words with continuous
exponents, supplies a single profinite word
for every compatible family.

At every finite three-primary level the image
of $N_B$ in the new factor is a three-group:
its image in $F/H$ is cyclic of order three
and its intersection with the geometric
kernel is a finite three-group.  Together
with $\mathscr D$ its image is finite nilpotent
of primes two and three, generated by the
$97$ fixed generators and mapping onto
$(\mathbb Z/6)^{97}$ through $G_6$.
The retained $G_6$ shadow is an isomorphism
on both Sylow Frattini quotients, so the
corrected map is onto this actual common
$N_B$ image.  It is then onto the common
source.  With the two-adic boundary factor,
take a lower-central index larger than the
nilpotence classes of all the finite
three-primary images and the old Ferrari
image.  E.41's exact stable two-primary
lower-central subgroup is supplied as a
pure factor.  The quotient is the nilpotent
common image just considered.  This proves
joint generation with the full E.41 tower
at every finite level.  E.42's actual
coprime $H$-image product gives the remaining
full pro-$p$ factors and joint generation.

For E.43 take cofinal old common sources
containing $P_{20}$.  Their new $H$ images
are solvable, and $B/H$ is solvable.  Their
actual common quotient with $G_L$ is therefore
exactly $S=G_L/C$, with the specified
$P_{20}\to S$ map, and the actual joint
source is $G_{\rm old}\times_SG_L$.
The surjection $\rho_L(R_H)=C$ of E.43
supplies a correction in
$R_H=\bigcap_jH^{[j]}$ matching the Galois
word in $G_L$.  It preserves every new
class-three image.  The exact common-quotient
argument gives generation of this fibre
product, as in E.44.15.
The further maps are the proved maps
\[
 \begin{gathered}
 G_L\longrightarrow G_A\xrightarrow{\psi}G_5,
 \qquad \psi\rho_A=\rho_5,
 \\
 j\longmapsto\bigl(j,\psi\pi_A\pi_L(j)\bigr).
 \end{gathered}
 \tag{E.47.25}
\]
of E.46.26.  They identify the old joint
source with its actual joint image after
adjoining $G_5$.  The same corrected word
therefore retains the five-point Hurwitz
value and the full common four-strand
target of E.46.27.

Precisely, \textup{(E.47.2)} parametrizes the
images under the actual map
$\rho:\widehat F'\to\mathscr D\times\mathcal B_3$
with its prescribed first coordinate.
For each such image and any one word $f$
realizing it, the set of actual profinite
words realizing it is exactly
\[
                    f\,\ker(\rho|_{\widehat F'}).
 \tag{E.47.26}
\]
This follows by multiplying two preimages
and is the complete comparison between
correction-image fibres and word fibres.
Adjoining the simultaneous target maps
replaces the kernel in this same formula
by their joint kernel.

The eight new programs in \texttt{code} have
shared prefix \texttt{bc\_fermat\_class\_three\_}
and suffixes \texttt{integral.py},
\texttt{norm\_smith\_probe.py},
\texttt{rational\_obstruction.py},
\texttt{curve\_obstruction.py},
\texttt{integral\_tests.py},
\texttt{graph\_tests.py}, \texttt{fibres.py}, and
\texttt{finite\_image.py}.  Their executed
certificates retain the tensor/Hall comparison,
the actual source word preimages, the Smith
matrix and saturation operations, every
polynomial row change, the central right
inverse, all nine residue systems, the full
$432$-point coverage, and the exact nonlifting
value.  All have status \texttt{PASS}.
\end{proof}
